\begin{abstract}
针对六区域异构 AI 任务与新能源、储能及电网协同调度问题，本文建立以任务执行区域、开工时刻和储能动作为决策，以运行成本、碳排放、网络时延、新能源利用率、服务质量及峰值净购电为评价指标，并受算力、功率、时延、期限、能量平衡和 SOC 边界约束的多目标模型。实际求解采用固定 Huber 滞后回归、确定性贪心任务调度、规则型 SOC 前向策略及固定种子情景复算，各指标经无量纲化后用于候选权衡，所得结果仅为通过硬约束审计的启发式可行方案。结果表明，50000 条任务全部通过校验，缺失数量为 0；预测总体 MAE 和 RMSE 分别为 24.973340521042196 GPU、35.69449198489753 GPU。最后 24 小时调度 538 个实际任务，其中 68 个在结清时域完成；碳感知调度产生 209 个迁移任务。固定负荷下，含储能运行成本为 -419516103.81021965 元，较无储能基线增加 304223.92909026146 元，碳排放、峰值净购电功率和绝对爬坡量变化均为 0.0，表明该规则策略未获得经济性与削峰收益。
\end{abstract}

\noindent\textbf{关键词：}异构 AI 任务；算力调度；Huber 回归；储能协同；碳感知调度；多目标权衡
